\itemhead{1.1 例题 · 例题2（教材p.3）}{XJ-B1-C01-EX-P003-Y1223}
\label{XJ-B1-C01-EX-P003-Y1223}
{\footnotesize 湘教版教材p.3 / PDF第10页 / 1.1 例题 / 例题2}
\field{原题}{下列集合中哪些是有限集？哪些是无限集？
\subq{1}一元二次方程 $x^2+1=0$ 的全体实根组成的集合 $A$；
\subq{2}所有素数（即质数）组成的集合 $B$；
\subq{3}所有矩形组成的集合 $C$；
\subq{4}所有亚洲国家组成的集合 $D$。}
\field{参考答案}{$A=\varnothing$，是有限集；$B,C$ 是无限集；$D$ 是有限集。}
{\small \field{关联知识}{S01 集合、元素与常用数集（教材帮p.1、2 / PDF10、11）}
\field{方法与易错点}{先辨认研究对象和数集范围，再代入定义判断；含参集合逐个来源分类。 易错：把自然数集误作正整数集；混淆元素和单元素集合；候选参数造成元素重复。}
\field{答案依据}{例题原解位置：湘教版教材p.3及题库原图；独立编写可编辑参考解答。答案册此部分主要收练习、习题，不伪造例题答案册位置。}}
